\section{Ethical and Broader Impacts}

\paragraph{Real World Impacts} Advancing the capabilities of autonomous agents comes with many broader considerations and ethical implications. Strong autonomous agents have the potential to improve the accessibility of computer-based tasks, potentially aiding individuals with disabilities or those lacking technical skills. More broadly, agents have the potential to automate large portions of routine computer work. While the capabilities of existing autonomous agents are insufficient for even simple tasks (as shown in this paper), these impacts highlight the need to ensure that the broader economic and social implications on employment are carefully considered if/when autonomous agents are deployed in real world applications.

\paragraph{Bias and Safety}  When developing autonomous agents, it is also imperative to ensure that these agents do not inadvertently exclude or disadvantage any group. Further analysis is essential to ensure that deployed agents do not exhibit unintended biases. Agents also have the potential to cause more harm (than regular LLMs) in real world applications if careful safeguards are not in place. Further research is necessary to understand and mitigate possible harmful behaviors.

\paragraph{Intended Uses} \BenchmarkName{} is a research benchmark to measure and evaluate the progress of multimodal agents. It is primarily meant to act as a self-contained sandbox environment for safely building robust agents. The models we presented in this paper are research prototypes, and not intended for deployment in practical applications in their current state (especially in high risk domains).